\sscp{**kindoR/**kVndoR(a)}{14-03-04}
Forms relating to **kindoR or **kandoRa (or similar) in \trf{14-06}
have a medial voiced plosive rather than a voiceless fricative
like **kusay, *kusor or {\#}kisor(o).
Only two have a medial nasal-consonant sequence /nd/.
(In the Table below,
STB = \tsc{Seram-Tanimbar-Bomberai}.)

\begin{table}[h]\centering\tcp{Forms resembling **kVndoR(a)}{14-06}
\begin{adjustbox}{width=\textwidth}
	\begin{threeparttable}\st{3pt}
	\begin{tabular}{l|lllll}\lsptoprule
Family	&	Subgroup	&	Language	&	ISO	&	form	&	gloss	\\	\midrule
\tsc{Austronesian}	&	STB, \tsc{E. Islands}	&	Geser	&	ges	&	kidor	&	cuscus	\\
\tsc{Austronesian}	&	STB, \tsc{E. Islands}	&	Gorom	&	ges	&	iˈdoʔ, iˈdo	&	cuscus	\\
\tsc{Austronesian}	&	STB, \tsc{E. Islands}	&	Kowiai	&	kwh	&	indor	&	cuscus	\\
\tsc{Austronesian}	&	STB, \tsc{E. Islands}	&	Kesui	&	(wah)	&	udora	&	cuscus	\\
\tsc{Austronesian}	&	STB, \tsc{E. Islands}	&	Watubela	&	wah	&	kadola	&	cuscus	\\
\tsc{Austronesian}	&	SHWNG	&	Yerisiam	&	ire	&	átóòrà	&	cuscus	\\
\tsc{Austronesian}	&	SHWNG	&	Matbat	&	xmt	&	doː (?), iba¹²ŋ 	&	cuscus	\\
\tsc{Austronesian}	&	SHWNG	&	As	&	asz	&	do	&	cuscus	\\
\tsc{Austronesian}	&	SHWNG	&	Buli	&	bzq	&	do	&	cuscus	\\
\tsc{Austronesian}	&	SHWNG	&	Gane	&	gzn	&	do	&	cuscus	\\
\tsc{Hatam-Mansim}	&		&	Hatam	&	had	&	indzɛya	&	opossum	\\
Isolate	&		&	Mor	&	moq	&	ˈiror, iora	&	tree wallaby	\\
\tsc{Tor-Orya}	&		&	Orya\su{†}	&	ury	&	sindor	&	marsupial	\\
Isolate	&		&	Mawes\su{†}	&	mgk	&	kudur	&	marsupial	\\
\tsc{Kwerba}	&		&	Isirawa\su{†}	&	srl	&	iˈroro	&	marsupial	\\
\lspbottomrule
	\end{tabular}
		\begin{tablenotes}\tnsize
			\item[†] {Orya, Mawes, and Isirawa are located further east than \frf{14-03} extends.}
		\end{tablenotes}
	\end{threeparttable}
\end{adjustbox}
\end{table}

It is possible that Buli
and Gane/Giman \it{do} could derive from a form like \mbox{**kandoRa},
since *k > {\0},
antepenultimate V > {\0},
*nd > \it{d}, *R > {\0} (working from \citealp{Kamholz2014}).
But they could equally derive from a form like **kindoR,
which collectively accounts for more languages in a wider distribution
in this immediate region than does **kandoRa.
Furthermore Gorom \it{iˈdo {\tl} iˈdoʔ} ‘cuscus’ has
final stress, which, when found in \tsc{Seram-Tanimbar-Bomberai} languages,
often signals which syllable will likely be preserved in SHWNG
languages that reduce to one syllable.

The initial high vowels varying between \it{i}
and \it{u} are not too unusual.
Synchronic variants with \it{u {\tl} i} are common cross-linguistically,
and are also frequent in \tsc{Seram-Tanimbar-Bomberai} languages,
along with some nearby \tsc{Ambon-Seram} languages (notably Manusela and Liana-Seti).
For example, Liana-Seti \it{\tbr{u}lafa} ‘rat’
and neighbouring Masiwang \it{\tbr{i}laha} ‘rat’ (another trisyllable
referring to another rodent-like animal; cf.
PMP *balabaw ‘rat, mouse,
rodent’); Manusela \it{wol\tbr{u}-ni} ‘gall,
bile’, Geser \it{fol\tbr{i}} ‘gall, bile’,
Watubela \it{fel\tbr{i}-n} ‘gall, bile’,
Kei Kecil \it{fir\tbr{u}-n} ‘gall, bile’, (cf.
PMP *qapəju ‘gall, bile’); Manusela \it{tol\tbr{u}-a} ‘egg’,
Masiwang \it{tol\tbr{i}-n} ‘egg’, (cf. PMP *qatəluR ‘egg’).
Examples with \it{u {\tl} i} variations in this region
and in these languages are not hard to find.

The correspondence of \it{r -- r -- r -- l} in Geser \it{kido\tbr{r}},
Kowiai \it{indo\tbr{r}}, Kesui \it{udo\tbr{r}a},
Watubela \it{kado\tbr{l}a} is regular,
as *R > \it{r} in Geser,
Kowiai, and Kesui, while *R > \it{l} in Watubela (\srf{11-03}).
These consonants are all as easily accounted for by **kindoR
as they are by **kandoRa.
Note also that Kesui *k > {\0} is regular.

But comparing Geser \it{kidor},
Kowiai \it{indor}, Kesui \it{udora},
Watubela \it{kadola}, (all Austronesian \tsc{Seram-Tanimbar-Bomberai}
languages of the \tsc{Eastern Islands} node) one must ask: did
Geser and Kowiai lose a final V? Or did Watubela
and Kesui add a final V? Since Geser
and Kowiai normally retain final a{\#}
and final r{\#}, there is an unexplained mismatch in syllable
structure here with Watubela.
Asserting the loss of the final historical *a in Geser,
Gorom, and Kowiai would be unexpected.
Consider trisyllables: Geser \it{tiliŋa} ‘ear’,
Kowiai \it{teriga} ‘ear’, Watubela \it{telŋa} ‘ear’ (< PMP *taliŋa
‘ear’); Geser \it{urana} ‘male’,
Kowiai \it{muana, muruana} ‘male’,
Watubela \it{maŋkana} ‘male’ (< PMP *maRuqanay ‘male’); Geser
\it{wawina} ‘female’, Kowiai \it{maɸina} ‘female’ (< PMP *ba-b<in>ahi,
*ma-b<in>ahi ‘female’); Kowiai \it{ɸituʔa} ‘intestines’ (< PMP
*bituka ‘stomach, large intestine’).
Final r{\#} is also preserved from some historical trisyllables,
as in Kowiai \it{asaʔar} ‘cassowary’; Geser \it{tolor} ‘egg’,
but not Kowiai \it{toro-n} ‘egg’,
Watubela \it{katlu} (< PMP *qatəluR ‘egg’)
and from historical disyllables: Geser \it{timor} ‘east monsoon’,
Kowiai \it{timor} ‘east monsoon’ (< PMP *timuR ‘south or east
wind’); Geser \it{laar} ‘sail’ (< PMP *layaR ‘sail’); Geser \it{fusal}
‘navel’, Kowiai \it{ɸusar} ‘navel’ (< PMP *pusəj ‘navel’); Kowiai \it{waʔar}
‘base’ (< PMP *wakaR ‘root’).
It would not be unreasonable from these data to conclude that
the final vowel of Watubela \it{kadola} is an addition,
if Geser \it{kidor}, Gorom \it{iˈdo {\tl} iˈdoʔ},
Kowiai \it{indor}, and Watubela \it{kadola} are from the same
source, particularly given that they all belong to the same node of the
same subgroup.\footnote{
		Geser often adds epenthetic final \it{a{\#}} or \it{ra{\#}}.
		For example, for the Waru variety of Geser on the Seram mainland
		it is stated (in \citeauthor{Stokhof1982a} [\citeyear[55]{Stokhof1982a}];
		stressed syllables are underlined)
		as “at the end of nouns one usually hears a soft \und{a} sound which
		is very weak, e.g.: \und{ko}tàn and also \und{ko}tàn(à), sīs\und{ī}
		and also sīs\und{ī}(à), \und{wà}tàn and also \und{wà}tàn(à).
		Many nouns which end in a vowel have a variant ending in a very
		weak r or rà,
		e.g.: īm\und{à} and also īm\und{ar} or īm\und{àr}(à),
		kat\und{oe}à and also kat\und{oe}àr or kat\und{oe}r(à),
		\und{a}fī and also \und{a}fīr or \und{a}fīr(à),
		b\und{ōī} and also b\und{ōī}r or b\und{ōī}r(à).” \citet{LoskiLoski1989} treat Geser
		and Gorom as two dialects of the Geser-Gorom language.
		For Gorom, \citet{Visser2019a} lists \it{iˈdoʔ} ‘cuscus’,
		and \citet[51]{Kakerissa1986} list \it{iˈdo} ‘cuscus’,
		both with final stress.
		Nearby languages such as Kei Kecil have ultimate stress as a rule.
		Geser and Watubela are classified within the same node together
		in all known sources,
		so it is not unreasonable to query the status of the final vowel
		of the Watubela trisyllable.}

Correspondences such as Geser \it{tolor} ‘egg’,
(< PMP *qatəluR ‘egg’); Geser \it{timor} ‘east monsoon’,
Kowiai \it{timor} ‘east monsoon’ (< PMP *timuR ‘south or east
wind’) also show that Geser
and Kowiai can derive \it{or{\#}} from a historical *uR,
and do not require a historical mid-vowel *o.
Nonetheless, *o does better reflect the on-the-ground data.\footnote{
		The Sera dialect of Fordata has a form /kandrua/ [kandruə] ‘rat’.
		(Schwa is an allophone of /a/ in open final syllables.) It is
		unclear whether this form is a reflex of PCEMP **kanzupay ‘rat’
		or of PCEMP **kandoRa.
		If it is from **kandoRa,
		that would further support a reconstruction with *u,
		**kanduRa, which also accounts for much of the regional data.}

In \srf{01-08} it was mentioned for the target region that “there
is a tendency towards a weakening of the pre-penultimate vowel
in historical trisyllabic roots.” Ante-penultimate vowel reduction
is a well attested phenomenon in the region,
and is also found in the languages in question,
as in PMP *maRuqanay ‘male’ > Geser \it{\tbr{u}rana} ‘male’,
Kowiai \it{m\tbr{u}ana, mur\tbr{u}ana} ‘male’,
but Watubela \it{m\tbr{a}ŋkana} ‘male’.
Once again, it is the Watubela form that is the odd man out among
related languages within the same node of the same subgroup.
Given regular sound correspondences,
the first vowel in Geser \it{kidor},
Gorom \it{iˈdo {\tl} iˈdoʔ},
Kowiai \it{indor}, and perhaps Kesui \it{udora},
are not explained by **kandoRa.
Given the range of forms in \trf{14-06},
what is unusual is the antepenultimate \it{a} found in Watubela \it{kadola}.
A correspondences of \it{u : a}
and \it{i : a} is so unusual,
that in relation to the variations in the antepenultimate vowel
in Watubela, Kowiai, Geser, and Kesui, \citet[268,
footnote 6]{Blust2012b} asserts that “\it{u : a}
and \it{i : a} are not attested vowel correspondences in any
languages of eastern Indonesia.”\footnote{
		In \srf{09-01} (\trf{09-05}, \prf{tab:09-05})
		we provide data that show this claim is incorrect.
		Similarly, most \tsc{Babar} langauegs have undergone
		and *a > \it{u} sound change (\trf{04-55}, \prf{tab:04-55}).}
While \citeauthor{Blust2012b}
is assuming that Watubela is the conservative form (from the
perspective of the Oceanic data),
within the context of similar forms in surrounding languages
in \trf{14-06}, it appears that Watubela antepenultimate \it{a} is the anomaly,
and not the other way around.
Determining the form of a reconstruction for eastern Indonesian
data based on the Oceanic data further away from the source is methodologically unsound.
We note that Geser,
Gorom, Kowiai, Watubela and Kesui all belong to the same \tsc{Eastern
Islands} node of the \tsc{Seram-Tanimbar-Bomberai} subgroup (see \crf{ch:STB}).

In and near our target region,
only Watubela and Yerisiam support the antepenultimate *a found in **kandoRa.
Methodologically, \tsc{POc} *kadoRa ‘cuscus’ represents only
one data point in EMP,
and PEMP represents only one data point in PCEMP (if those nodes are accepted).
But given that the full form of **kandoRa is fully supported
by only one (!) language outside of EMP,
and that this form is anomalous when compared to related forms
in nearby languages, it is unsound to determine the form of a higher level node by
the form of one lower level node that is in conflict with forms
in other nodes that theoretically link in at an earlier level.\footnote{
		We also note in passing that several of the Oceanic forms (Nauna
		\it{kocay}, Penchal \it{kotay}, Sori \it{ohay},
		Seimat \it{koxa}) appear on the surface to more closely resemble
		Sulawesi **kusay than the \tsc{POc} *kadoRa form to which they
		are ascribed \citep{BlustTrussel2020}.
		These forms are also flagged as being anomalous in their own
		region, appealing to metathesis to explain their vowels with regards
		to \tsc{POc} *kadoRa \citep[241]{Blust1982}.}

A further mystery that remains to be solved is the wider distribution
of reflexes of **kVndoR(a) or purported PCEMP **kandoRa ‘cuscus’.
There are two pockets of languages separated by about 1,900 kilometers
that appear to reflect these forms (see map in \citealp[62]{Schapper2011}).
The data in the western pocket (the one in our target region)
is better accounted for by **kindoR(a) or **kinduR(a) ‘cuscus’.
The eastern pocket appears to be adequately accounted for by
\tsc{POc} *kadoRa \citep[241]{Blust1982},
or \tsc{POc} *k(ʷ)adroRa ‘cuscus’ \citep[225]{OsmondPawley2011}.
The close resemblance in both form
and meaning between these two distant pockets appears to be too
remarkable to be dismissed as chance or parallel development.
Is all the relevant data available? An explanation has remained elusive to date.
Nevertheless, the best way to explain the diverse reflexes of
**kVndoR(a) is through borrowing.
Given this, it would be premature to assign **kVndoR(a) to the
parent of the \tsc{Eastern Islands} node.